\section{Introduction}

Zeckendorf expansion singles out a canonical Fibonacci representation of every
positive integer.  The problem considered here is a finite-window version of
that principle.  For each $m\ge 1$ we consider the binary window space
\[
\Omega_m:=\{0,1\}^m.
\]
Let
\[
X_m:=\{w\in\{0,1\}^m:\text{$w$ contains no subword }11\},
\]
and define
\[
\Fold_m(\omega):=\pi_m(Z(N(\omega))),
\]
where $N(\omega)$ is the Fibonacci value of $\omega$, $Z(N(\omega))$ is its
Zeckendorf expansion, and $\pi_m$ records the first $m$ digits.  Thus
$\Fold_m$ replaces an arbitrary binary window by its canonical admissible
representative in Fibonacci numeration.

Two questions motivate the paper.  The first is local: can $\Fold_m$ be
described by a finite normalization procedure, and how does it behave when the
window length changes?  The second is reconstructive: if one slides this
normalization along a bi-infinite binary sequence, how much of the original
sequence can be recovered from the resulting list of overlapping folded
windows?  In number-theoretic terms, the second question asks whether the
vector of overlapping local Fibonacci weighted sums modulo $F_{m+2}$
determines an arbitrary binary block uniquely, and if so, what the sharp
observation length is.

We answer the first question by showing that $\Fold_m$ is the normal-form map
of a finite terminating confluent rewrite system; see
Proposition~\ref{prop:rewrite-newman} and
Corollary~\ref{cor:fold-order-independent}.  We also prove that the maps
$\Fold_m$ are compatible with restriction in the window length.  In particular,
the finite admissible languages form an inverse system whose limit is the
one-sided golden-mean shift; see Theorem~\ref{thm:inverse-limit-golden}.

To study the second question, for each $m$ we consider the sliding block code
\[
\Phi_m:\{0,1\}^{\ZZ}\to X_m^{\ZZ}
\]
with image $Y_m$, the language of overlapping normalized windows.  The main
result is that for every $m\ge 3$ and every $n\ge m$, each admissible block of
$n$ consecutive folded symbols has a unique raw lift of length $n+m-1$.  Thus
the threshold $m$ is exactly the point at which local normalization becomes
lossless on finite blocks.  This yields the block count
\[
\lvert\mathcal L_n(Y_m)\rvert=2^{n+m-1}
\qquad (n\ge m),
\]
and gives a family of local inverse rules obtained by reading any prescribed
coordinate of the unique length-$m$ lift.  The same block bijection also
describes the Bernoulli measures pushed forward from the full two-shift; see
Theorem~\ref{thm:block-bijection-threshold},
Corollary~\ref{cor:block-counts-threshold},
Corollary~\ref{cor:higher-block-identification}, and
Theorems~\ref{thm:Phi-m-conjugacy-threshold},
\ref{thm:bernoulli-transport-cylinder-law}, and
\ref{thm:exact-markov-order-bernoulli}.  At the one-sided endpoint this
recovers an inverse with memory $m-1$, and no smaller one-sided memory
suffices; see Corollary~\ref{cor:Phi-m-memory-threshold-exact}.

Thus the same integer $m$ governs several different aspects of the problem.  It
is the first scale at which overlapping folded windows determine the
underlying raw block, the Markov memory of the push-forward of Bernoulli
measures, the minimal synchronizing length of the natural presentation, and
the size parameter of the right Fischer cover.  The symbolic-dynamical
language serves mainly to record these consequences efficiently.

\subsection*{Related work}

The arithmetic background begins with Zeckendorf's theorem
\cite{Zeckendorf1972}.  Finite-state aspects of numeration systems were
developed systematically by Frougny and collaborators
\cite{Frougny1991FibonacciFiniteAutomata,Frougny1992NumbersAutomata,
Frougny1999OnlineAddition,FrougnySolomyak1992}, while algorithmic forms of
Zeckendorf arithmetic, including carry propagation and parallel normalization,
were studied in \cite{AhlbachUsatineFrougnyPippenger2013}.  For broader
numeration settings one may compare Ostrowski systems and their
automata-theoretic realizations
\cite{Berthe2001Ostrowski,HieronymiTerry2018OstrowskiAddition,
BaranwalaSchaefferShallit2021OstrowskiAutomatic}, as well as more general work
on non-standard and linear-recurrence numeration systems with finite-memory
automaton structure and normalization algorithms
\cite{AlloucheShallit2003AutomaticSequences,
FrougnyPelantovaSvobodova2013MinimalDigitSets,
Frougny2012NumbersAsStreamsLeiden,LegerskySvobodova2018EWM}.  Recent work has
also connected Zeckendorf representations with sequence-level statistical and
dynamical phenomena; see \cite{ManiboMiroRustTadeo2020ZeckendorfMixing}.

The distributional and combinatorial structure of Zeckendorf decompositions has
been investigated by Grabner, Tichy, and collaborators from the viewpoint of
digital problems and asymptotic distribution of digit sums in non-standard
numeration systems \cite{GrabnerTichy1995FibonacciDigitalSums}; by Miller and
collaborators, who established Gaussian behavior for the number of summands in
Zeckendorf representations \cite{MillerWang2012GaussianZeckendorf}; and from the
ergodic-theoretic perspective connecting beta-expansions and Parry conditions to
symbolic dynamics of numeration \cite{Parry1960BetaExpansions,
Renyi1957Representations}.

These strands of literature address different questions: recognition of legal
expansions, arithmetic operations in non-standard numeration systems,
automaticity, distributional behavior, or digit-sum statistics.  The present
paper does not study addition, parallel arithmetic, or asymptotic digit
statistics.  Instead it fixes a window length $m$, normalizes every window
locally in Fibonacci numeration, and asks when the original binary block can be
reconstructed exactly from the resulting local Zeckendorf data.  Equivalently,
the maps $\Phi_m$ form a family of finite-window factor maps attached to
Fibonacci numeration, and the later finite-state constructions record the
corresponding determinization data.

More concretely, the work of Frougny and collaborators shows how finite automata
recognize admissible expansions and implement normalization or arithmetic in
Fibonacci and related numeration systems.  Our question is different.  We begin
with an arbitrary binary block, apply truncated Zeckendorf normalization on
every overlapping window of length $m$, and ask when those local normal forms
still determine the original block exactly.  Unlike the above literature, the
present paper proves a sharp local congruence rigidity threshold and derives its
exact consequences for representation counting and random local residue
distributions.  We are not aware of an earlier result, even in the Fibonacci
setting, that identifies this sharp observation length for the recoverability
problem or the corresponding threshold in the natural deterministic
presentation.

This threshold is of arithmetic interest because it measures how much
information survives finite truncation of Zeckendorf normalization.  The answer
is not merely qualitative.  Once the window length reaches $m$, every
admissible folded block of length $n\ge m$ corresponds to exactly one binary
block of length $n+m-1$, so
$\lvert\mathcal L_n(Y_m)\rvert=2^{n+m-1}$.  The same bijection gives explicit
push-forward formulas for Bernoulli measures and an exact $m$-step Markov law
on the folded alphabet.

In purely number-theoretic terms, the main results can be stated as follows.
For each binary block $\omega=\omega_1\cdots\omega_L\in\{0,1\}^L$ with
$L=n+m-1$, the $t$-th local Fibonacci sum is
$S_t^{(m)}(\omega):=\sum_{j=0}^{m-1}\omega_{t+j}F_{j+2}$, and its residue
modulo $F_{m+2}$ is
$R_t^{(m)}(\omega):=S_t^{(m)}(\omega)\bmod F_{m+2}$.  The residue vector
$\mathcal R_{m,n}(\omega):=(R_1^{(m)}(\omega),\ldots,R_n^{(m)}(\omega))$
determines $\omega$ uniquely.  The symbolic-dynamical language provides a
convenient framework for organizing the proofs of this and the following
consequences.

\begin{theorem}[Main results]
\label{thm:intro-main}
For every $m\ge 3$, the following statements hold.
\begin{enumerate}[label=\textup{(\roman*)}]
\item\label{item:main-congruence-rigidity}
\textup{(Local Fibonacci congruence rigidity.)}
For every $n\ge m$ and $L=n+m-1$, the map
$\mathcal R_{m,n}:\{0,1\}^L\to\{0,1,\ldots,F_{m+2}-1\}^n$ that sends a binary
block to its vector of overlapping local Fibonacci sums modulo $F_{m+2}$ is
injective.  Equivalently, the congruence system
\[
\sum_{j=0}^{m-1}\omega_{t+j}F_{j+2}\equiv r_t\pmod{F_{m+2}}
\qquad (1\le t\le n)
\]
has at most one solution $\omega\in\{0,1\}^L$ for each
$(r_1,\ldots,r_n)\in\{0,1,\ldots,F_{m+2}-1\}^n$.
\item\label{item:main-block-bijection}
\textup{(Exact block reconstruction.)}
For every $n\ge m$, every admissible block of $n$ consecutive folded
symbols has a unique raw lift of length $n+m-1$.  Equivalently, the block map
$B_{m,n}:\{0,1\}^{n+m-1}\to\mathcal L_n(Y_m)$ is a bijection.
Consequently
$\lvert\mathcal L_n(Y_m)\rvert=2^{n+m-1}$ for every $n\ge m$, and the
$n$-block presentation of the normalized language is canonically the
$(n+m-1)$-block presentation of the full binary shift.
\item\label{item:main-bit-recovery}
\textup{(Exact bit recovery from residue windows.)}
Write $\mathfrak{R}_{m,m}$ for the set of realizable length-$m$ residue
vectors.  For each $0\le s\le m-1$, there exists a function
$\rho_{m,s}:\mathfrak{R}_{m,m}\to\{0,1\}$ such that
\[
\omega_{t+s}=\rho_{m,s}\bigl(R_t^{(m)}(\omega),\ldots,R_{t+m-1}^{(m)}(\omega)\bigr)
\]
for every $\omega\in\{0,1\}^{n+m-1}$ and every valid index $t$.  The minimum
number of consecutive residues required for any one-sided recovery rule is
exactly $m$.
\item\label{item:main-bernoulli}
\textup{(Explicit law and exact Markov order.)}
For every $p\in(0,1)$, if $\nu_p$ denotes the Bernoulli product law on
$\{0,1\}^{\ZZ}$ and $\mu_{m,p}:=(\Phi_m)_*\nu_p$, then every cylinder
$[a]_0\subset Y_m$ with $a\in\mathcal L_n(Y_m)$ and $n\ge m$ satisfies the
explicit formula
\[
\mu_{m,p}([a]_0)=p^{|\omega(a)|_1}(1-p)^{n+m-1-|\omega(a)|_1},
\qquad
\omega(a):=B_{m,n}^{-1}(a).
\]
Hence
$H_{\mu_{m,p}}(Y_0^{n-1})=(n+m-1)h_2(p)$, the excess entropy equals
$(m-1)h_2(p)$, and the coordinate process under $\mu_{m,p}$ is exact
$m$-step Markov.  The random local Fibonacci residue process
$R_t:=(\sum_{j=0}^{m-1}X_{t+j}F_{j+2})\bmod F_{m+2}$ under i.i.d.\
Bernoulli$(p)$ input is itself an exact $m$-step Markov process.
\item\label{item:main-inverse}
\textup{(Sliding block inverses.)}
Exact recovery from overlapping normalized windows is possible from any
length-$m$ observation window: for each $0\le r\le m-1$ there is a sliding
block inverse with memory $r$ and anticipation $m-1-r$.  At the one-sided
endpoint $r=m-1$, no smaller memory suffices.
\item\label{item:main-fischer}
\textup{(Fischer cover and presentation complexity.)}
Over the natural alphabet $X_m$, the canonical minimal right-resolving
presentation of $Y_m$ is the explicit suffix graph $G_m$, so it has exactly
$2^{m-1}$ states and $2^m$ edges.  Since $\lvert X_m\rvert=F_{m+2}$, the
natural deterministic presentation size grows exponentially relative to the
alphabet size.
\end{enumerate}
\end{theorem}

Taken together, these results show that a single integer controls the
finite-window arithmetic of the problem: it is the point at which local
Zeckendorf normalization becomes lossless on blocks, the sharp threshold for
local Fibonacci congruence rigidity, the exact Markov memory of the random
residue process, and the scale parameter of the natural deterministic
presentation; see again
Corollary~\ref{cor:local-complexity-scale-coincidence} and
Theorem~\ref{thm:natural-state-complexity-gap}.

These statements are proved in
Theorem~\ref{thm:block-bijection-threshold},
Corollary~\ref{cor:block-counts-threshold},
Corollary~\ref{cor:higher-block-identification},
Theorem~\ref{thm:Phi-m-conjugacy-threshold},
Theorem~\ref{thm:bernoulli-transport-cylinder-law},
Corollary~\ref{cor:exact-entropy-profile},
Theorem~\ref{thm:exact-markov-order-bernoulli},
Corollary~\ref{cor:Phi-m-memory-threshold-exact},
Theorem~\ref{thm:Y-m-exact-SFT-order},
Corollary~\ref{cor:uniform-synchronizing-length},
Proposition~\ref{prop:ambiguity-shell-nilpotency}, and
Corollary~\ref{cor:G-m-fischer-cover}, together with
Theorem~\ref{thm:natural-state-complexity-gap}.

\paragraph{Sketch of the reconstruction mechanism.}
The inverse is built in two steps.  First, the pairwise overlap maps
$\Theta_m$ compress adjacent folded symbols to the fold of their common raw
overlap; iterating this reduction through a block
$a_1\cdots a_m\in\mathcal L_m(Y_m)$ yields the central bit
$\Theta_{m,m}(a_1\cdots a_m)$.  Second, once this central bit and the folded
inner overlaps are known, the one-step left- and right-resolving properties of
$G_m$ force the neighboring bits uniquely, so an outward induction recovers the
entire raw lift of length $2m-1$.  The inverse family in Section~7 simply
records the same reconstruction after choosing which coordinate of the
recovered length-$m$ raw window is read as the output bit.

The paper is organized as follows.  Sections~2 and~3 develop finite-window
normalization and its rewrite description.  Section~4 proves compatibility in
the window length, Section~5 formulates the overlap-decoding problem,
Section~6 records the explicit finite-state presentation used later, and
Section~7 proves the block reconstruction theorem, derives the measure-theoretic
consequences, and then studies the natural right-resolving presentation of the
folded language.  Section~8 reformulates the main results in purely
number-theoretic terms: local Fibonacci congruence rigidity, representation
fingerprint separation, exact bit recovery from residue windows, and the
precise distribution and Markov order of random local residue vectors.

